\subsection{齐性空间上的拟不变测度}

引理 3. --- 设 G 是局部紧群，\(\mu\) 是 G 上的左 Haar 测度，\(\nu\) 和 \(\nu'\) 是 G 上的两个非零拟不变测度。若对于每个 \(s \in G\)，\(\boldsymbol{\gamma}(s)\nu\) 相对于 \(\nu\) 的密度与 \(\boldsymbol{\gamma}(s)\nu'\) 相对于 \(\nu'\) 的密度局部 \(\mu\)-几乎处处相等，则 \(\nu\) 与 \(\nu'\) 成比例。

写成 \(\nu = \rho \cdot \mu\)、\(\nu' = \rho' \cdot \mu\)，其中 \(\rho, \rho'\) 是局部 \(\mu\)-可积且处处非零的函数，定义在 \(G\) 上（§1, No.~9, Prop. 11）。对于每个 \(s \in G\)，

\[
\boldsymbol {\gamma} (s) \nu = (\boldsymbol {\gamma} (s) \rho) \cdot \mu , \quad \boldsymbol {\gamma} (s) \nu^ {\prime} = (\boldsymbol {\gamma} (s) \rho^ {\prime}) \cdot \mu ,
\]

且假设蕴含 \(\rho^{-1} \cdot \pmb{\gamma}(s) \rho = \rho'^{-1} \cdot \pmb{\gamma}(s) \rho'\) 局部 \(\mu\)-几乎处处成立。令 \(\sigma = \rho'/\rho\)，这是一个 \(\mu\)-可测函数，定义在 \(G\) 上。对于每个 \(s \in G\)，\(\pmb{\gamma}(s) \sigma = \sigma\) 局部 \(\mu\)-几乎处处成立。因此 \(\sigma\) 局部 \(\mu\)-几乎处处等于一个常数，由将 \(X = H = G\) 应用于 Prop. 6 的 Cor. 2。

设 G 是局部紧群，H 是 G 的闭子群。考虑关于 H 的左陪集齐性空间 G/H，G 在其上连续地从左作用。我们将证明，G/H 上非零拟不变测度只有一个等价类。

注意，H 通过右平移在 G 上连续且 逆紧地作用；而商空间正是 G/H，是仿紧的（GT, III, §4, No.~6, Prop. 13）。因此可在 X = G 的情形应用 No. 1 至 No. 4 的结果。于是有映射 \(f \mapsto f^{b}\)，从 \(\mathcal{H}(G)\) 到 \(\mathcal{K}(G/H)\)，以及映射 \(\lambda \mapsto \lambda^{\sharp}\)，从 \(\mathcal{M}(G/H)\) 到 \(\mathcal{M}(G)\)（一旦固定 H 上的左 Haar 测度 \(\beta\)）。G 在 G/H 上从左作用这一事实还给出如下附加性质：

\[
\boldsymbol {\gamma} _ {\mathrm{G/H}} (s) \cdot f ^ {b} = \left(\boldsymbol {\gamma} _ {\mathrm{G}} (s) \cdot f\right) ^ {b} \quad \left(s \in \mathrm{G}, f \in \mathscr {K} (\mathrm{G})\right)\tag{15}
\]

\[
\left(\boldsymbol {\gamma} _ {\mathrm{G/H}} (s) \cdot \lambda\right) ^ {\sharp} = \boldsymbol {\gamma} _ {\mathrm{G}} (s) \cdot \lambda^ {\sharp} \quad \left(s \in \mathrm{G}, \lambda \in \mathscr {M} (\mathrm{G/H})\right).\tag{16}
\]

事实上，对于每个 \(x \in G\)，

\[
\begin{array}{r l} & {\left(\boldsymbol {\gamma} _ {\mathrm{G/H}} (s) \cdot f ^ {b}\right) (\pi (x)) = f ^ {b} \big (s ^ {- 1} \pi (x) \big) = f ^ {b} \big (\pi (s ^ {- 1} x) \big)} \\ & {\qquad = \int_ {\mathrm{H}} f (s ^ {- 1} x \xi) d \beta (\xi) = \int_ {\mathrm{H}} \big (\boldsymbol {\gamma} _ {\mathrm{G}} (s) f \big) (x \xi) d \beta (\xi) = \big (\boldsymbol {\gamma} _ {\mathrm{G}} (s) f \big) ^ {b} \big (\pi (x) \big),} \end{array}
\]

从而得到公式 (15)，进而推出公式 (16)。

引理 4. --- 设 \(\lambda\) 是 G/H 上的测度 \(\neq 0\)，\(\mu\) 是 G 上的左 Haar 测度。以下性质等价：

\begin{enumerate}
\def\labelenumi{\alph{enumi})}
\item
  \(\lambda\) 关于 G 拟不变；
\item
  G/H 的子集 A 局部 \(\lambda\)-可忽略的充要条件是，\(\overline{\pi}^{1}(A)\) 局部 \(\mu\)-可忽略；
\item
  测度 \(\lambda^{\sharp}\) 与 \(\mu\) 等价。
\end{enumerate}

假设这些性质成立，并令 \(\lambda^{\sharp} = \rho \cdot \mu\)，其中 \(\rho\) 是局部 \(\mu\)-可积且处处非零的函数。则对于每个 \(s \in \mathbf{G}\)，其密度 \(\theta_s\) 是 \(\gamma_{\mathrm{G / H}}(s)\lambda\) 相对于 \(\lambda\) 的密度，满足

\[
\theta_ {s} (\pi (x)) = \frac {\rho (s ^ {- 1} x)}{\rho (x)}\tag{17}
\]

在 G 上局部 \(\mu\)-几乎处处成立。

\begin{enumerate}
\def\labelenumi{\alph{enumi})}
\setcounter{enumi}{2}
\item
  \(\Rightarrow\) b)：这立即由 Prop. 6 a) 得出。
\item
  \(\Rightarrow\) a)：若性质 b) 成立，则 G/H 的局部 \(\lambda\)-可忽略子集所成的集合关于 G 不变，故 \(\lambda\) 关于 G 拟不变。
\item
  \(\Rightarrow\) c)：假设 \(\lambda\) 关于 G 拟不变；对于每个 \(s\in G\)，\(\lambda\) 与 \(\pmb{\gamma}_{\mathrm{G / H}}(s)\lambda\) 等价，因此 \(\lambda^{\sharp}\) 与 \(\pmb{\gamma}_{\mathrm{G}}(s)\cdot \lambda^{\sharp} = (\pmb{\gamma}_{\mathrm{G / H}}(s)\cdot \lambda)^{\sharp}\) 等价（Prop. 6 的 Cor. 1）；由于 \(\lambda^{\sharp}\neq 0\)，由 §1, No.~9, Prop. 11，\(\lambda^{\sharp}\) 与 \(\mu\) 等价。
\end{enumerate}

此外，对于每个 \(s \in \mathbf{G}\)，

\[
\begin{array}{r l} & {(\theta_ {s} \circ \pi) \cdot \lambda^ {\sharp} = (\theta_ {s} \cdot \lambda) ^ {\sharp} = (\boldsymbol {\gamma} _ {\mathrm{G/H}} (s) \lambda) ^ {\sharp} = \boldsymbol {\gamma} _ {\mathrm{G}} (s) \lambda^ {\sharp}} \\ & {\qquad \qquad \qquad \qquad \qquad = (\boldsymbol {\gamma} _ {\mathrm{G}} (s) \rho) \cdot \mu = \frac {\boldsymbol {\gamma} _ {\mathrm{G}} (s) \rho}{\rho} \cdot \lambda^ {\sharp},} \end{array}
\]

从而得到 (17)。

a) 与 b) 的等价性首先蕴含前面所述的唯一性结论，甚至给出一个更精确的结果：

定理 1. --- 设 G 是局部紧群，H 是 G 的闭子群。

\begin{enumerate}
\def\labelenumi{\alph{enumi})}
\item
  G/H 上任意两个非零拟不变测度都等价；对于这些测度，局部可忽略的 G/H 子集正是其在 G 中的逆像对于某个 Haar 测度局部可忽略的那些子集。
\item
  设 \(\lambda, \lambda'\) 是 G/H 上的两个非零拟不变测度。若对于每个 \(s \in \mathbf{G}\)，\(\gamma_{\mathrm{G / H}}(s)\lambda\) 相对于 \(\lambda\) 的密度与 \(\gamma_{\mathrm{G / H}}(s)\lambda'\) 相对于 \(\lambda'\) 的密度对于 \(\lambda\)（或 \(\lambda'\)）几乎处处相等，则 \(\lambda\) 与 \(\lambda'\) 成比例。
\end{enumerate}

断言 a) 立即由引理 4 得出。设 \(\lambda\) 和 \(\lambda'\) 是满足 b) 条件的两个非零拟不变测度。则对于每个 \(s \in G\)，\(\gamma_G(s)\lambda^\sharp\) 相对于 \(\lambda^\sharp\) 的密度与 \(\gamma_G(s)\lambda'^\sharp\) 相对于 \(\lambda'^\sharp\) 的密度局部 \(\mu\)-几乎处处相等，因此由引理 3，\(\lambda^\sharp\) 与 \(\lambda'^\sharp\) 成比例，从而 \(\lambda\) 与 \(\lambda'\) 成比例。

另一方面，引理 4 将寻找 G/H 上的非零拟不变测度归结为寻找 G 上与 Haar 测度等价且形如

\(\lambda^{\sharp}\) 的测度。在此问题上有如下引理：

引理 5. --- 设 \(\mu\) 是 G 上的左 Haar 测度，\(\rho\) 是局部 \(\mu\)-可积函数。\(\rho \cdot \mu\) 具有 \(\lambda^{\sharp}\) 的形式的充要条件是，对于每个 \(\xi \in \mathbf{H}\)，

\[
\rho (x \xi) = \frac {\Delta_ {\mathrm{H}} (\xi)}{\Delta_ {\mathrm{G}} (\xi)} \rho (x)\tag{18}
\]

在 G 上局部 \(\mu\)-几乎处处成立。

说 \(\rho \cdot \mu\) 具有 \(\lambda^{\sharp}\) 的形式，等价于说对于每个 \(\xi \in \mathbf{H}\)，有 \(\delta(\xi)(\rho \cdot \mu) = \Delta_{\mathrm{H}}(\xi)\rho \cdot \mu\)（Prop. 4）。现在，

\[
\boldsymbol {\delta} (\xi) (\rho \cdot \mu) = (\boldsymbol {\delta} (\xi) \rho) \cdot (\boldsymbol {\delta} (\xi) \mu) = \Delta_ {G} (\xi) (\boldsymbol {\delta} (\xi) \rho) \cdot \mu ,
\]

引理得证。

现在可以建立前面所述的存在性结论，甚至得到一个更精确的结果：

定理 2. --- 设 G 是局部紧群，H 是 G 的闭子群，μ 是 G 上的左 Haar 测度，β 是 H 上的左 Haar 测度。

\begin{enumerate}
\def\labelenumi{\alph{enumi})}
\tightlist
\item
  存在连续的函数 \(\rho\)，取值为 \(>0\)，定义在 \(\mathbf{G}\) 上，使得
\end{enumerate}

\[
\rho (x \xi) = \frac {\Delta_ {\mathrm{H}} (\xi)}{\Delta_ {\mathrm{G}} (\xi)} \rho (x)
\]

对所有 \(x \in \mathbf{G}\) 和 \(\xi \in \mathbf{H}\) 都成立。

\begin{enumerate}
\def\labelenumi{\alph{enumi})}
\setcounter{enumi}{1}
\item
  给定这样的函数 \(\rho\)，可构造测度 \(\lambda = (\rho \cdot \mu) / \beta\)，定义在 \(\mathbf{G} / \mathbf{H}\) 上，且 \(\lambda\) 是关于 \(\mathbf{G}\) 拟不变的非零正测度。
\item
  对于 \(s, x\) 属于 \(\mathbf{G}\)，\(\rho(sx) / \rho(x)\) 只依赖于 \(s\) 和 \(\pi(x)\)，因而定义了函数 \(\chi\)，连续且取值为 \(>0\)，定义在 \(\mathbf{G} \times (\mathbf{G} / \mathbf{H})\) 上，使得
\end{enumerate}

\[
\chi (s, \pi (x)) = \frac {\rho (s x)}{\rho (x)}.\tag{19}
\]

则

\[
\boldsymbol {\gamma} _ {\mathrm{G/H}} (s) \lambda = \chi (s ^ {- 1},.) \cdot \lambda \quad \text {   for   all   } s \in \mathrm{G}.\tag{20}
\]

\begin{enumerate}
\def\labelenumi{\alph{enumi})}
\item
  由 Prop. 7 得出。
\item
  由引理 5 和引理 4 得出。
\item
  由 (17) 得出。
\end{enumerate}

注记。--- 1) 由 No.~3 的注记 1 可推出，G/H 上的非零拟不变测度正是 G 上某个 Haar 测度在 \(\pi\) 下的拟像测度。

*2) 若 G 是 Lie 群，稍后将看到，Th. 2 中的函数 \(\rho\) 可以选为无穷可微函数。*

在 Th. 2 的条件下，结合 Ch. V, §4, Th. 2 和 Prop. 2 将关于 \(\mu\) 的性质转化为关于 \(\rho \cdot \mu\) 的性质，No. 3 和 No. 4 的某些结果可具体表述如下：

\begin{enumerate}
\def\labelenumi{\alph{enumi})}
\item
  设 \(f\) 是一个 \(\mu\)-可测函数，定义在 \(G\) 上，取值于一个拓扑空间，并且在某个 \(\mu\)-可积集的可数并之外为常值；则那些 \(\dot{x} \in G / H\) 使函数 \(\xi \mapsto f(x\xi)\) 不是 \(\beta\)-可测的点所成的集合是局部 \(\lambda\)-可忽略的。
\item
  设 \(f\) 是一个 \(\mu\)-可测的 \(\geqslant 0\) 函数，定义在 \(G\) 上，并在某个 \(\mu\)-可积集的可数并之外为零。则函数
\end{enumerate}

\[
\dot {x} \mapsto \int_ {\mathrm{H}} ^ {*} f (x \xi) d \beta (\xi)
\]

在 G/H 上是 λ-可测的，并且

\[
\int_ {\mathrm{G}} ^ {*} f (x) \rho (x) d \mu (x) = \int_ {\mathrm{G} / \mathrm{H}} ^ {*} d \lambda (\dot {x}) \int_ {\mathrm{H}} ^ {*} f (x \xi) d \beta (\xi) \quad (\dot {x} = \pi (x)).
\]

\begin{enumerate}
\def\labelenumi{\alph{enumi})}
\setcounter{enumi}{2}
\tightlist
\item
  设 \(\mathbf{f}\) 是一个 \(\rho \cdot \mu\)-可积函数，定义在 \(\mathbf{G}\) 上，取值于 Banach 空间或 \(\overline{\mathbf{R}}\)。则那些 \(\dot{x} \in \mathrm{G} / \mathrm{H}\) 使函数 \(\xi \mapsto \mathbf{f}(x\xi)\) 不是 \(\beta\)-可积的点所成的集合是 \(\lambda\)-可忽略的；函数 \(\dot{x} \mapsto \int_{\mathrm{H}} \mathbf{f}(x\xi) d\beta(\xi)\) 是 \(\lambda\)-可积的，并且
\end{enumerate}

\[
\int_ {\mathrm{G}} \mathbf {f} (x) \rho (x) d \mu (x) = \int_ {\mathrm{G/H}} d \lambda (\dot {x}) \int_ {\mathrm{H}} \mathbf {f} (x \xi) d \beta (\xi).
\]

\begin{enumerate}
\def\labelenumi{\alph{enumi})}
\setcounter{enumi}{3}
\tightlist
\item
  存在连续函数 \(h \geqslant 0\)，定义在 \(G\) 上，其支集与 \(KH\) 的交为紧，其中 \(K\) 是 \(G\) 的紧子集，并且 \(\int_{H} h(x\xi) d\beta(\xi) = 1\) 对每个 \(x \in G\) 成立。对于函数 \(g\)，定义在 \(G/H\) 上，关于 \(\lambda\) 可测（分别为局部可积、本质可积、可积）的充要条件是 \(h \cdot (g \circ \pi)\) 关于 \(\rho \cdot \mu = \lambda^{\sharp}\) 具有相应性质；并且，当 \(g\) 关于 \(\lambda\) 本质可积时，有
\end{enumerate}

\[
\int_ {\mathrm{G/H}} \mathbf {g} (u) d \lambda (u) = \int_ {\mathrm{G}} h (x) \mathbf {g} (\pi (x)) \rho (x) d \mu (x).
\]

